% =====================================================================
\chapter{Συμπεράσματα και Μελλοντική Εργασία}
\label{ch:conclusions}
% =====================================================================

\section{Σύνοψη}

Η παρούσα εργασία παρουσίασε το \Probolos, μια πύλη αποδοχής συσκευών \en{USB}
με προεπιλεγμένη άρνηση για το \en{Linux}. Το σύστημα αξιοποιεί τη διεπαφή
εξουσιοδότησης του πυρήνα ώστε κάθε νέα συσκευή να εμφανίζεται απενεργοποιημένη,
την εξετάζει σε έως τέσσερα στάδια (ταυτότητα, συνέπεια, συμπεριφορά,
αποθήκευση) και αφήνει την τελική απόφαση στον άνθρωπο, με τριβή ανάλογη του
κινδύνου. Ένα ιστορικό με κανονικοποιημένα αποτυπώματα επιτρέπει τον εντοπισμό
συσκευών που άλλαξαν μετά την έγκρισή τους, ενώ ο διαχωρισμός προνομίων
περιορίζει τον κώδικα που εκτελείται ως \en{root} σε ένα σαφώς ορισμένο σύνολο
έξι αιτημάτων.

\section{Κύρια συμπεράσματα}

\begin{enumerate}
  \item \textbf{Η προεπιλεγμένη άρνηση είναι εφικτή σε αμετάβλητο πυρήνα.} Ο
        μηχανισμός \code{authorized\_default} του \en{Linux} αρκεί για να
        αποτραπεί η πρόσδεση οδηγών σε άγνωστες συσκευές, χωρίς τροποποιήσεις
        στον πυρήνα ή εικονικοποίηση.
  \item \textbf{Οι έλεγχοι συνέπειας εντοπίζουν τη «φθηνή» κατηγορία
        επιθέσεων.} Η πλειονότητα των εμπορικά διαθέσιμων συσκευών \en{BadUSB}
        δηλώνει συνδυασμούς λειτουργιών χωρίς νόμιμη χρήση· ένας απλός κανόνας
        όπως ο \code{storage-with-keyboard} τους εντοπίζει σε μικροδευτερόλεπτα.
  \item \textbf{Ο ρυθμός πληκτρολόγησης δεν είναι αξιόπιστο κριτήριο, η
        σιωπή είναι.} Η προσομοίωση έδειξε ότι ο κανόνας ρυθμού παρακάμπτεται
        με ελάχιστη τυχαιότητα, ενώ ο κανόνας «πληκτρολόγησε χωρίς να το
        αγγίζει κανείς» παραμένει ισχυρός για κάθε ωφέλιμο φορτίο που εκτελείται
        κατά την καραντίνα.
  \item \textbf{Η ασφάλεια έναντι κλειδώματος είναι πρωταρχική απαίτηση.} Τα
        σοβαρότερα ελαττώματα που βρέθηκαν (P0) αφορούσαν ακριβώς τις διαδρομές
        επαναφοράς της πύλης. Ένα εργαλείο ασφαλείας που κλειδώνει τον χρήστη
        έξω θα απεγκατασταθεί — και τότε δεν προστατεύει κανέναν.
  \item \textbf{Ο συστηματικός έλεγχος κώδικα και οι αυτοματοποιημένες τεχνικές
        είναι συμπληρωματικά.} Ο πρώτος εντόπισε τα περισσότερα και σοβαρότερα
        ελαττώματα· οι δεύτερες βρήκαν ελαττώματα που κανένας άνθρωπος δεν είχε
        εντοπίσει.
\end{enumerate}

\section{Μελλοντική εργασία}
\label{sec:future}

\begin{itemize}
  \item \textbf{Ανάλυση περιγραφέων αναφορών \en{HID}.} Ο αναλυτής υπάρχει ήδη
        (συνάρτηση \code{walk\_hid\_items} της μονάδας \code{descriptors\_safe}), αλλά ο περιγραφέας αναφορών
        δεν είναι διαθέσιμος στο \en{sysfs} όσο η συσκευή είναι απενεργοποιημένη.
        Η ανάγνωσή του κατά την καραντίνα ή με μεταφορά ελέγχου (\en{control
        transfer}) θα επέτρεπε νέο ισχυρό κανόνα: συσκευή που δηλώνει «ποντίκι»
        ενώ ο περιγραφέας αναφορών της περιγράφει πληκτρολόγιο.
  \item \textbf{Παρακολούθηση μετά την έγκριση με \en{HID-BPF}}, ώστε να
        ανιχνεύεται «υπομονετική» συσκευή που επιτίθεται αργότερα.
  \item \textbf{Ενεργή ανάκριση ως ανιχνευτής.} Το ερευνητικό εργαλείο
        \code{interrogate.py} συλλέγει δεδομένα για το πώς απαντούν συσκευές σε
        ασυνήθιστα αιτήματα· αν τα δεδομένα διαχωρίζουν συνηθισμένες συσκευές από
        μικροελεγκτές γενικής χρήσης, θα μπορούσε να γίνει κανόνας.
  \item \textbf{\en{HMAC} στον χώρο εμπιστοσύνης}, με κλειδί προσβάσιμο μόνο από
        τον \en{root}.
  \item \textbf{Επεξήγηση της μετάλλαξης}: τι ακριβώς άλλαξε μεταξύ δύο
        εμφανίσεων μιας συσκευής και αν η αλλαγή είναι εύλογη (π.χ. ενημέρωση
        υλικολογισμικού).
  \item \textbf{Αληθοφάνεια κατασκευαστή}: \en{VID} του οποίου οι συμβολοσειρές
        κατονομάζουν άλλη εταιρεία.
  \item \textbf{Εμπειρική μελέτη ψευδώς θετικών} σε μεγάλο δείγμα πραγματικών
        συσκευών, ξεκινώντας από τους κανόνες που προέρχονται από την
        προδιαγραφή και όχι από παρατηρήσεις (\code{descriptor-chain-truncated},
        \code{descriptor-length-overstated}).
  \item \textbf{Μελέτη χρηστικότητας} με χρήστες, για τη μέτρηση της
        αποτελεσματικότητας της κλιμακούμενης τριβής.
\end{itemize}

\section{Επίλογος}

Το \en{USB} σχεδιάστηκε για να «απλώς λειτουργεί». Η εργασία αυτή υποστηρίζει ότι
ένα σύστημα μπορεί να διατηρήσει αυτή την ευκολία για τις νόμιμες συσκευές,
ζητώντας από τον χρήστη ένα ή δύο κλικ, ενώ ταυτόχρονα αφαιρεί από μια κακόβουλη
συσκευή το βασικό της πλεονέκτημα: την αυτόματη εμπιστοσύνη του λειτουργικού
συστήματος. Όπως υποδηλώνει και το όνομά του — \emph{πρόβολος}, το εμπόδιο που
τοποθετείται μπροστά — το \Probolos{} δεν υπόσχεται απόλυτη προστασία· υπόσχεται
ότι τίποτα δεν περνά χωρίς να το δει πρώτα κάποιος.
